%=============================================================
%  第3章 データと分析手法
%=============================================================
\chapter{データと分析手法}
\label{chap:data-method}

\section{データ}
\label{sec:data}

\subsection{サンプルと期間}

本研究では，日本上場企業 \num{3558} 社の財務データを用いる。
サンプル期間は 2010 年度から 2024 年度であり，
延べ \num{49642} 社・年の観測値を含む（表\ref{tab:data-summary}参照）。

データは〇〇（データベース名）から取得した。
サンプルから除外した観測値として，(1) 財務諸表が入手できないもの，
(2) 分析に用いる変数のすべてが欠損しているものが挙げられる。

\begin{table}[htbp]
  \centering
  \caption{サンプルの基本統計量}
  \label{tab:data-summary}
  % TODO: Python でテーブルを生成して \input{tables/tab_summary_stats.tex} に置き換える
  \begin{tabular}{lrrrr}
    \toprule
    変数 & 観測数 & 平均 & 標準偏差 & 中央値 \\
    \midrule
    総資産（百万円） & \num{49642} & \num{285688} & \num{1812459} & \num{27882} \\
    売上高（百万円） & \num{49642} & \num{190391} & \num{852424} & \num{26887} \\
    自己資本比率（\%） & \num{49642} & 51.8 & 21.0 & 52.3 \\
    \bottomrule
  \end{tabular}
\end{table}

\subsection{変数の定義}

分析に用いる財務変数は表\ref{tab:variables}に示す 11 変数である。

\begin{table}[htbp]
  \centering
  \caption{分析変数の定義}
  \label{tab:variables}
  \begin{tabularx}{\linewidth}{llX}
    \toprule
    変数名 & 区分 & 定義・出典 \\
    \midrule
    総資産     & 規模  & 貸借対照表の資産合計（百万円） \\
    売上高     & 規模  & 損益計算書の売上高（百万円） \\
    短期借入金 & 負債  & 流動負債のうち短期借入金（百万円），欠損値は 0 を代入 \\
    長期借入金 & 負債  & 固定負債のうち長期借入金（百万円），欠損値は 0 を代入 \\
    利益剰余金 & 内部留保 & 純資産のうち利益剰余金（百万円） \\
    現金同等物 & 流動性 & キャッシュ・フロー計算書の期末残高（百万円） \\
    営業利益   & 収益性 & 損益計算書の営業利益（百万円） \\
    営業 CF    & CF    & キャッシュ・フロー計算書の営業活動によるキャッシュ・フロー \\
    投資 CF    & CF    & 投資活動によるキャッシュ・フロー \\
    財務 CF    & CF    & 財務活動によるキャッシュ・フロー \\
    自己資本比率 & 安全性 & 純資産合計 ÷ 資産合計 $\times$ 100（\%） \\
    \bottomrule
  \end{tabularx}
\end{table}

\section{データの前処理}
\label{sec:preprocessing}

\subsection{対数変換}

財務変数は企業規模によって大きく分布が歪んでいる。
通常の対数変換 $\log(1+x)$ は $x < -1$ の場合に欠損値を生成し，
赤字企業や負の利益剰余金を持つ企業をサンプルから除外する
選択バイアスをもたらす。

本研究では，この問題に対処するため，以下の符号付き対数変換を採用する。

\begin{equation}
  \mathrm{signed\_log1p}(x) = \mathrm{sign}(x) \cdot \log(1 + |x|)
  \label{eq:signed-log1p}
\end{equation}

この変換は，正の値に対しては通常の $\log(1+x)$ と同一であり，
負の値に対しては符号を保持しつつ大きさを圧縮する。
これにより，損失企業や財務的に脆弱な企業をサンプルに含めることができ，
選択バイアスを回避できる。

\subsection{欠損値の処理}

短期借入金と長期借入金については，有価証券報告書を持つ企業が
借入を行っていない場合に欠損として計上されるケースが多い。
本研究では，このような構造的欠損に対して 0 を代入する。

\section{クラスタリング手法}
\label{sec:clustering}

\subsection{年度内標準化}

各財務変数は，年度ごとに平均 0・標準偏差 1 に標準化する。

\begin{equation}
  \tilde{x}_{i,t} = \frac{x_{i,t} - \bar{x}_t}{s_t}
  \label{eq:zscore}
\end{equation}

ここで $\bar{x}_t$ と $s_t$ は年度 $t$ における当該変数の標本平均と標本標準偏差である。
この年度内標準化により，各変数が年度をまたいで比較可能になる。

\subsection{主成分分析（PCA）}

年度内標準化後の財務変数行列に対して主成分分析（PCA）を適用し，
多重共線性を解消した上で次元削減を行う。

本研究では，予測モデルの訓練期間と一致させるために，
訓練期間（2010〜2021年度）のデータのみを用いて PCA を推定し，
全期間にわたるデータを同一の主成分基底に射影する。

表\ref{tab:pca-variance}に各主成分の寄与率を示す。
第1〜第3主成分の累積寄与率は 64.5\% であり，
本研究ではこれら3つの主成分スコアをクラスタリングの入力変数として用いる。

\begin{table}[htbp]
  \centering
  \caption{主成分の寄与率}
  \label{tab:pca-variance}
  % TODO: \input{tables/tab_pca_variance.tex} に置き換える
  \begin{tabular}{crr}
    \toprule
    主成分 & 寄与率（\%） & 累積寄与率（\%） \\
    \midrule
    PC1 & 35.0 & 35.0 \\
    PC2 & 19.5 & 54.5 \\
    PC3 & 10.0 & 64.5 \\
    \bottomrule
  \end{tabular}
\end{table}

\subsection{K-Means クラスタリングとハンガリアン法によるラベル整合}

本研究では，年度ごとに独立した K-Means クラスタリングを実施した上で，
ハンガリアン法（Hungarian algorithm）によって年度間のクラスターラベルを整合させる。

年度 $t$ の K-Means クラスタリングにより，各企業を $K=5$ 個のクラスターに分類する。
クラスター数 $K$ の選択には Silhouette 係数とエルボー法を用いた。

年度をまたいだラベルの整合は，前年度の重心と当年度の重心の間のユークリッド距離行列に
対して線形割当問題（ハンガリアン法）を解くことで実現する。

\begin{equation}
  \hat{\sigma}_t = \arg\min_{\sigma \in S_K}
  \sum_{k=1}^{K} \left\| \mathbf{c}_{t-1, k} - \mathbf{c}_{t, \sigma(k)} \right\|^2
  \label{eq:hungarian}
\end{equation}

ここで $\mathbf{c}_{t,k}$ は年度 $t$ のクラスター $k$ の重心ベクトルであり，
$S_K$ は $K$ 個の要素の置換全体の集合である。

\section{遷移分析}
\label{sec:transition}

\subsection{遷移行列の推定}

年度 $t$ でクラスター $j$ に属する企業が，年度 $t+1$ でクラスター $k$ に
遷移する確率を遷移行列の要素 $p_{jk}$ として推定する。

\subsection{二値プロビット（遷移の有無）}

翌年度に異なるクラスターへ遷移するかどうかを被説明変数とする
二値プロビットモデルを推定する。

\begin{equation}
  \Pr(y_{i,t} = 1) = \Phi(\mathbf{x}_{i,t}' \boldsymbol{\beta})
  \label{eq:probit}
\end{equation}

説明変数 $\mathbf{x}_{i,t}$ には，(1) 主成分スコア（PC1〜PC3），
(2) 対数変換済み財務変数，(3) 現在クラスターダミー，
(4) 年度固定効果が含まれる。

\subsection{多項ロジット（遷移先の決定）}

遷移が発生した企業を対象に，遷移先クラスターを被説明変数とする
多項ロジットモデルを推定する。

\section{予測モデル}
\label{sec:prediction}

\subsection{XGBoost}

本研究では，勾配ブースティング木（Gradient Boosted Trees）の実装である
XGBoost \citep{ChenGuestrin2016} を予測モデルとして用いる。

\paragraph{Step 1: 遷移の二値分類}

遷移するか否かを目的変数とする二値分類モデルを推定する。
非遷移（約 80\%）と遷移（約 20\%）の不均衡に対処するため，
\texttt{scale\_pos\_weight} パラメータを非遷移数／遷移数の比率（3.99）に設定する。

時系列バイアスを防ぐため，訓練期間を 2010〜2021 年度，
テスト期間を 2022〜2024 年度とする時系列分割を採用する。

\paragraph{Step 2: 遷移先の多クラス分類}

遷移した企業のみを対象に，遷移先クラスター（0〜4）を目的変数とする
多クラス分類モデル（\texttt{objective: multi:softprob}）を推定する。

\subsection{SHAP による解釈可能性分析}

各特徴量の予測への貢献度を Shapley Additive Explanations（SHAP；
\tcite{LundbergLee2017}）によって定量化する。
本研究では XGBoost 内蔵の TreeSHAP を使用し，
各特徴量のMean $|\mathrm{SHAP}|$ を重要度指標として用いる。

\section{ロバストネスチェック}
\label{sec:robustness}

分析の設計上の懸念点（PCA における情報漏洩の可能性）に対して，
以下の3つのロバストネスチェックを実施する。

\begin{description}
  \item[方法①（Train-only クラスタリング）]
    PCA の推定を訓練期間（2010〜2021 年度）のデータに限定し，
    テスト期間はその主成分基底で変換する。これにより，
    テスト期間の分散構造が PCA 成分の計算に混入することを防ぐ。

  \item[方法②（ロバストネス比較）]
    全期間 PCA（Pattern A）と Train-only PCA（Pattern B）の
    両クラスタリングで予測モデルを推定し，性能を比較する。
    AUC の差分が十分に小さければ，PCA のリークの影響が限定的であることを示す。

  \item[方法③（クラスター変数除外）]
    現在クラスターを特徴量から除外したモデルを推定することで，
    財務変数単独の説明力を検証する。
\end{description}
